% Generated. Do not edit.
% Projection scriptural-date-location/section, version 2: the Scriptural Date and
% Location page, generated whole from research/chronology.toml, the Scripture
% chronology corpus and the source library by tools/tpt document-types refresh.
% document: liturgy/roman-rite/1962/propers/temporal/01-first-sunday-of-advent
\runninghead{Scriptural Date and Location}
\section{Scriptural Date and Location}
\label{sec:date-location}

{\small Each passage the Mass reads or sings from Scripture, in canonical order. The Context column distinguishes the event or setting, an attributed author's era, and the writing or formation of the text. The Date column gives calendar orientation. A proposal span summarizes the extent of numerical dates in the sources, including gaps between proposals. Undated occasions and events remain distinct from composition dates. The complete alternatives and their evidence follow the table.\par}
\medskip
% Generated. Do not edit.
% Deterministic projection of research/chronology.toml.
% schema: 5
% document: liturgy/roman-rite/1962/propers/temporal/01-first-sunday-of-advent
% generated-by: tools/tpt proper-chronology annotations
\newcommand{\chronologyannotationclaim}[6]{#6}
\newcommand{\chronologyannotationreach}[3]{}
\newcommand{\chronologyannotationgroup}[3]{#3}
\newcommand{\chronologyannotation}[1]{%
  \ifcsname triptychchronologyannotation@#1\endcsname
    \csname triptychchronologyannotation@#1\endcsname
  \else
    \PackageError{triptych}{No chronology annotation for #1}{}%
  \fi}
% introit: Introit
\expandafter\def\csname triptychchronologyannotation@introit\endcsname{%
  \chronologyannotationgroup{traditional-attribution}{disputed}{\textbf{Traditional attribution} -- disputed: \chronologyannotationreach{Ps.24.1}{Ps.24}{true}\chronologyannotationreach{Ps.24.2}{Ps.24}{true}\chronologyannotationreach{Ps.24.3}{Ps.24}{true}\chronologyannotationreach{Ps.24.4}{Ps.24}{true}\chronologyannotationclaim{israel.monarchy.david-traditional-era}{traditional-attribution}{catholic-traditional-v1}{disputed}{reigned from 1055 to 1015 B.C.}{David (reign in the usual chronology), B.C. 1055--1015}.}%
  \space \chronologyannotationgroup{historical-setting}{research-pending}{\textbf{Historical setting} -- Occasion date unresolved.}%
  \space \chronologyannotationgroup{composition}{preferred}{\textbf{Composition}: \chronologyannotationreach{Ps.24.1}{Ps}{true}\chronologyannotationreach{Ps.24.2}{Ps}{true}\chronologyannotationreach{Ps.24.3}{Ps}{true}\chronologyannotationreach{Ps.24.4}{Ps}{true}\chronologyannotationclaim{critical.psalms.latest-composition-boundary}{composition}{catholic-critical-v1}{preferred}{before the Maccabean period, around 165 B.C.; no individual psalm can be dated securely}{Before c. 165 B.C.}}%
}
% epistle: Epistle
\expandafter\def\csname triptychchronologyannotation@epistle\endcsname{%
  \chronologyannotationgroup{composition}{disputed}{\textbf{Composition} -- disputed: \chronologyannotationreach{Rom.13.11}{Rom}{true}\chronologyannotationreach{Rom.13.12}{Rom}{true}\chronologyannotationreach{Rom.13.13}{Rom}{true}\chronologyannotationreach{Rom.13.14}{Rom}{true}\chronologyannotationclaim{composition.epistle-to-the-romans}{composition}{catholic-traditional-v1}{disputed}{(Romans), 57}{A.D. 57}; \chronologyannotationreach{Rom.13.11}{Rom}{true}\chronologyannotationreach{Rom.13.12}{Rom}{true}\chronologyannotationreach{Rom.13.13}{Rom}{true}\chronologyannotationreach{Rom.13.14}{Rom}{true}\chronologyannotationclaim{composition.epistle-to-the-romans}{composition}{catholic-traditional-v1}{disputed}{the Epistle to the Romans (58)}{A.D. 58}.}%
}
% gradual: Gradual
\expandafter\def\csname triptychchronologyannotation@gradual\endcsname{%
  \chronologyannotationgroup{traditional-attribution}{disputed}{\textbf{Traditional attribution} -- disputed: \chronologyannotationreach{Ps.24.3}{Ps.24}{true}\chronologyannotationreach{Ps.24.4}{Ps.24}{true}\chronologyannotationclaim{israel.monarchy.david-traditional-era}{traditional-attribution}{catholic-traditional-v1}{disputed}{reigned from 1055 to 1015 B.C.}{David (reign in the usual chronology), B.C. 1055--1015}.}%
  \space \chronologyannotationgroup{historical-setting}{research-pending}{\textbf{Historical setting} -- Occasion date unresolved.}%
  \space \chronologyannotationgroup{composition}{preferred}{\textbf{Composition}: \chronologyannotationreach{Ps.24.3}{Ps}{true}\chronologyannotationreach{Ps.24.4}{Ps}{true}\chronologyannotationclaim{critical.psalms.latest-composition-boundary}{composition}{catholic-critical-v1}{preferred}{before the Maccabean period, around 165 B.C.; no individual psalm can be dated securely}{Before c. 165 B.C.}}%
}
% alleluia: Alleluia
\expandafter\def\csname triptychchronologyannotation@alleluia\endcsname{%
  \chronologyannotationgroup{historical-setting}{research-pending}{\textbf{Historical setting} -- Occasion date unresolved.}%
  \space \chronologyannotationgroup{composition}{preferred}{\textbf{Composition}: \chronologyannotationreach{Ps.84.8}{Ps.41 Ps.42 Ps.43 Ps.44 Ps.45 Ps.46 Ps.47 Ps.48 Ps.83 Ps.84 Ps.86 Ps.87}{true}\chronologyannotationclaim{composition.psalms-of-the-sons-of-korah}{composition}{catholic-traditional-v1}{preferred}{between the days of Isaias and the return from exile}{Between the days of Isaias and the return from exile}.}%
}
% gospel: Gospel
\expandafter\def\csname triptychchronologyannotation@gospel\endcsname{%
  \chronologyannotationgroup{narrated-event}{preferred}{\textbf{Event}: \chronologyannotationreach{Luke.21.25}{Luke.21.5-36}{false}\chronologyannotationreach{Luke.21.26}{Luke.21.5-36}{false}\chronologyannotationreach{Luke.21.27}{Luke.21.5-36}{false}\chronologyannotationreach{Luke.21.28}{Luke.21.5-36}{false}\chronologyannotationreach{Luke.21.29}{Luke.21.5-36}{false}\chronologyannotationreach{Luke.21.30}{Luke.21.5-36}{false}\chronologyannotationreach{Luke.21.31}{Luke.21.5-36}{false}\chronologyannotationreach{Luke.21.32}{Luke.21.5-36}{false}\chronologyannotationreach{Luke.21.33}{Luke.21.5-36}{false}\chronologyannotationclaim{life-of-christ.prophecies-in-the-temple-lucan-account}{narrated-event}{catholic-traditional-v1}{preferred}{c. A.D. 29}{Jesus foretells the fall of Jerusalem and the coming of the Son of Man, as Luke places the discourse, c. A.D. 29 (derived)}.}%
  \space \chronologyannotationgroup{composition}{disputed}{\textbf{Composition} -- disputed: \chronologyannotationreach{Luke.21.25}{Luke}{true}\chronologyannotationreach{Luke.21.26}{Luke}{true}\chronologyannotationreach{Luke.21.27}{Luke}{true}\chronologyannotationreach{Luke.21.28}{Luke}{true}\chronologyannotationreach{Luke.21.29}{Luke}{true}\chronologyannotationreach{Luke.21.30}{Luke}{true}\chronologyannotationreach{Luke.21.31}{Luke}{true}\chronologyannotationreach{Luke.21.32}{Luke}{true}\chronologyannotationreach{Luke.21.33}{Luke}{true}\chronologyannotationclaim{composition.gospel-of-luke}{composition}{catholic-traditional-v1}{disputed}{About the year 70}{c. A.D. 70}; \chronologyannotationreach{Luke.21.25}{Luke}{true}\chronologyannotationreach{Luke.21.26}{Luke}{true}\chronologyannotationreach{Luke.21.27}{Luke}{true}\chronologyannotationreach{Luke.21.28}{Luke}{true}\chronologyannotationreach{Luke.21.29}{Luke}{true}\chronologyannotationreach{Luke.21.30}{Luke}{true}\chronologyannotationreach{Luke.21.31}{Luke}{true}\chronologyannotationreach{Luke.21.32}{Luke}{true}\chronologyannotationreach{Luke.21.33}{Luke}{true}\chronologyannotationclaim{composition.gospel-of-luke}{composition}{catholic-traditional-v1}{disputed}{before the end of the Roman imprisonment, when the Acts was finished}{Before the end of the Roman imprisonment, when the Acts was finished}.}%
}
% offertory: Offertory
\expandafter\def\csname triptychchronologyannotation@offertory\endcsname{%
  \chronologyannotationgroup{traditional-attribution}{disputed}{\textbf{Traditional attribution} -- disputed: \chronologyannotationreach{Ps.24.1}{Ps.24}{true}\chronologyannotationreach{Ps.24.2}{Ps.24}{true}\chronologyannotationreach{Ps.24.3}{Ps.24}{true}\chronologyannotationclaim{israel.monarchy.david-traditional-era}{traditional-attribution}{catholic-traditional-v1}{disputed}{reigned from 1055 to 1015 B.C.}{David (reign in the usual chronology), B.C. 1055--1015}.}%
  \space \chronologyannotationgroup{historical-setting}{research-pending}{\textbf{Historical setting} -- Occasion date unresolved.}%
  \space \chronologyannotationgroup{composition}{preferred}{\textbf{Composition}: \chronologyannotationreach{Ps.24.1}{Ps}{true}\chronologyannotationreach{Ps.24.2}{Ps}{true}\chronologyannotationreach{Ps.24.3}{Ps}{true}\chronologyannotationclaim{critical.psalms.latest-composition-boundary}{composition}{catholic-critical-v1}{preferred}{before the Maccabean period, around 165 B.C.; no individual psalm can be dated securely}{Before c. 165 B.C.}}%
}
% communion: Communion
\expandafter\def\csname triptychchronologyannotation@communion\endcsname{%
  \chronologyannotationgroup{historical-setting}{research-pending}{\textbf{Historical setting} -- Occasion date unresolved.}%
  \space \chronologyannotationgroup{composition}{preferred}{\textbf{Composition}: \chronologyannotationreach{Ps.84.13}{Ps.41 Ps.42 Ps.43 Ps.44 Ps.45 Ps.46 Ps.47 Ps.48 Ps.83 Ps.84 Ps.86 Ps.87}{true}\chronologyannotationclaim{composition.psalms-of-the-sons-of-korah}{composition}{catholic-traditional-v1}{preferred}{between the days of Isaias and the return from exile}{Between the days of Isaias and the return from exile}.}%
}
% presentation: orientation-v3
\newcommand{\chronologyorientationrow}[2]{%
  \ifcsname triptychchronologyorientation@#1\endcsname
    \csname triptychchronologyorientation@#1\endcsname{#2}%
  \else
    \PackageError{triptych}{No chronology orientation row for #1}{}%
  \fi}
\expandafter\def\csname triptychchronologyorientation@row-ffde50926b618362\endcsname#1{%
Introit: Ps. 24:1--3 (Vg.; Heb. 25), Ps. 24:4 (Vg.; Heb. 25) & Traditional attribution: David's reign in the usual chronology; Hebron, then Jerusalem & c.~1055--1015 B.C. (usual chronology; Corbett also reports dates 30--50 years later)\chronologyannotationreach{Ps.24.1}{Ps.24}{true}\chronologyannotationreach{Ps.24.2}{Ps.24}{true}\chronologyannotationreach{Ps.24.3}{Ps.24}{true}\chronologyannotationreach{Ps.24.4}{Ps.24}{true}\chronologyannotationclaim{israel.monarchy.david-traditional-era}{traditional-attribution}{catholic-traditional-v1}{disputed}{reigned from 1055 to 1015 B.C.}{}\\*[.1em]%
}
\expandafter\def\csname triptychchronologyorientation@row-1884ebeacdf710a8\endcsname#1{%
 & Historical setting & Occasion date unresolved\\*[.1em]%
}
\expandafter\def\csname triptychchronologyorientation@row-cdc5f95d42da0398\endcsname#1{%
 & Composition: Psalms & \chronologyannotationreach{Ps.24.1}{Ps}{true}\chronologyannotationreach{Ps.24.2}{Ps}{true}\chronologyannotationreach{Ps.24.3}{Ps}{true}\chronologyannotationreach{Ps.24.4}{Ps}{true}\chronologyannotationclaim{critical.psalms.latest-composition-boundary}{composition}{catholic-critical-v1}{preferred}{before the Maccabean period, around 165 B.C.; no individual psalm can be dated securely}{Before c.~165 B.C.}\\[.3em]%
}
\expandafter\def\csname triptychchronologyorientation@row-e043e4e3e3d8610c\endcsname#1{%
Offertory: Ps. 24:1--3 (Vg.; Heb. 25) & Traditional attribution: David's reign in the usual chronology; Hebron, then Jerusalem & c.~1055--1015 B.C. (usual chronology; Corbett also reports dates 30--50 years later)\chronologyannotationreach{Ps.24.1}{Ps.24}{true}\chronologyannotationreach{Ps.24.2}{Ps.24}{true}\chronologyannotationreach{Ps.24.3}{Ps.24}{true}\chronologyannotationclaim{israel.monarchy.david-traditional-era}{traditional-attribution}{catholic-traditional-v1}{disputed}{reigned from 1055 to 1015 B.C.}{}\\*[.1em]%
}
\expandafter\def\csname triptychchronologyorientation@row-59e074ad4fb4bf82\endcsname#1{%
 & Historical setting & Occasion date unresolved\\*[.1em]%
}
\expandafter\def\csname triptychchronologyorientation@row-8f3ee9c054245e05\endcsname#1{%
 & Composition: Psalms & \chronologyannotationreach{Ps.24.1}{Ps}{true}\chronologyannotationreach{Ps.24.2}{Ps}{true}\chronologyannotationreach{Ps.24.3}{Ps}{true}\chronologyannotationclaim{critical.psalms.latest-composition-boundary}{composition}{catholic-critical-v1}{preferred}{before the Maccabean period, around 165 B.C.; no individual psalm can be dated securely}{Before c.~165 B.C.}\\[.3em]%
}
\expandafter\def\csname triptychchronologyorientation@row-7d7d97f8358968c7\endcsname#1{%
Gradual: Ps. 24:3 (Vg.; Heb. 25), Ps. 24:4 (Vg.; Heb. 25) & Traditional attribution: David's reign in the usual chronology; Hebron, then Jerusalem & c.~1055--1015 B.C. (usual chronology; Corbett also reports dates 30--50 years later)\chronologyannotationreach{Ps.24.3}{Ps.24}{true}\chronologyannotationreach{Ps.24.4}{Ps.24}{true}\chronologyannotationclaim{israel.monarchy.david-traditional-era}{traditional-attribution}{catholic-traditional-v1}{disputed}{reigned from 1055 to 1015 B.C.}{}\\*[.1em]%
}
\expandafter\def\csname triptychchronologyorientation@row-f5e722ff79a9e57e\endcsname#1{%
 & Historical setting & Occasion date unresolved\\*[.1em]%
}
\expandafter\def\csname triptychchronologyorientation@row-6c4c7f6537f3262a\endcsname#1{%
 & Composition: Psalms & \chronologyannotationreach{Ps.24.3}{Ps}{true}\chronologyannotationreach{Ps.24.4}{Ps}{true}\chronologyannotationclaim{critical.psalms.latest-composition-boundary}{composition}{catholic-critical-v1}{preferred}{before the Maccabean period, around 165 B.C.; no individual psalm can be dated securely}{Before c.~165 B.C.}\\[.3em]%
}
\expandafter\def\csname triptychchronologyorientation@row-09b31c90e5cc8331\endcsname#1{%
Alleluia: Ps. 84:8 (Vg.; Heb. 85) & Historical setting & Occasion date unresolved\\*[.1em]%
}
\expandafter\def\csname triptychchronologyorientation@row-cf755db1bdb2f73a\endcsname#1{%
 & Composition: The Psalms of the Sons of Korah & c.~740--536 B.C. (Korahite collection estimate)\chronologyannotationreach{Ps.84.8}{Ps.41 Ps.42 Ps.43 Ps.44 Ps.45 Ps.46 Ps.47 Ps.48 Ps.83 Ps.84 Ps.86 Ps.87}{true}\chronologyannotationclaim{composition.psalms-of-the-sons-of-korah}{composition}{catholic-traditional-v1}{preferred}{between the days of Isaias and the return from exile}{}\\[.3em]%
}
\expandafter\def\csname triptychchronologyorientation@row-eb87f9559d582d9f\endcsname#1{%
Communion: Ps. 84:13 (Vg.; Heb. 85) & Historical setting & Occasion date unresolved\\*[.1em]%
}
\expandafter\def\csname triptychchronologyorientation@row-16c5e1487f2aee6d\endcsname#1{%
 & Composition: The Psalms of the Sons of Korah & c.~740--536 B.C. (Korahite collection estimate)\chronologyannotationreach{Ps.84.13}{Ps.41 Ps.42 Ps.43 Ps.44 Ps.45 Ps.46 Ps.47 Ps.48 Ps.83 Ps.84 Ps.86 Ps.87}{true}\chronologyannotationclaim{composition.psalms-of-the-sons-of-korah}{composition}{catholic-traditional-v1}{preferred}{between the days of Isaias and the return from exile}{}\\[.3em]%
}
\expandafter\def\csname triptychchronologyorientation@row-03a893478d4fede6\endcsname#1{%
Gospel: Luke 21:25--33 & Narrated event: Jesus foretells Jerusalem's fall and His coming; Jerusalem, by the Temple in Luke's narrative; Mount Olivet in Matthew and Mark & \chronologyannotationreach{Luke.21.25}{Luke.21.5-36}{false}\chronologyannotationreach{Luke.21.26}{Luke.21.5-36}{false}\chronologyannotationreach{Luke.21.27}{Luke.21.5-36}{false}\chronologyannotationreach{Luke.21.28}{Luke.21.5-36}{false}\chronologyannotationreach{Luke.21.29}{Luke.21.5-36}{false}\chronologyannotationreach{Luke.21.30}{Luke.21.5-36}{false}\chronologyannotationreach{Luke.21.31}{Luke.21.5-36}{false}\chronologyannotationreach{Luke.21.32}{Luke.21.5-36}{false}\chronologyannotationreach{Luke.21.33}{Luke.21.5-36}{false}\chronologyannotationclaim{life-of-christ.prophecies-in-the-temple-lucan-account}{narrated-event}{catholic-traditional-v1}{preferred}{c. A.D. 29}{c.~A.D.~29 (derived; Maas's chronology)}\\*[.1em]%
}
\expandafter\def\csname triptychchronologyorientation@row-ec61c0aa1cf6d38e\endcsname#1{%
 & Composition: The Gospel of St Luke & Disputed: \chronologyannotationreach{Luke.21.25}{Luke}{true}\chronologyannotationreach{Luke.21.26}{Luke}{true}\chronologyannotationreach{Luke.21.27}{Luke}{true}\chronologyannotationreach{Luke.21.28}{Luke}{true}\chronologyannotationreach{Luke.21.29}{Luke}{true}\chronologyannotationreach{Luke.21.30}{Luke}{true}\chronologyannotationreach{Luke.21.31}{Luke}{true}\chronologyannotationreach{Luke.21.32}{Luke}{true}\chronologyannotationreach{Luke.21.33}{Luke}{true}\chronologyannotationclaim{composition.gospel-of-luke}{composition}{catholic-traditional-v1}{disputed}{About the year 70}{c.~A.D.~70}; \chronologyannotationreach{Luke.21.25}{Luke}{true}\chronologyannotationreach{Luke.21.26}{Luke}{true}\chronologyannotationreach{Luke.21.27}{Luke}{true}\chronologyannotationreach{Luke.21.28}{Luke}{true}\chronologyannotationreach{Luke.21.29}{Luke}{true}\chronologyannotationreach{Luke.21.30}{Luke}{true}\chronologyannotationreach{Luke.21.31}{Luke}{true}\chronologyannotationreach{Luke.21.32}{Luke}{true}\chronologyannotationreach{Luke.21.33}{Luke}{true}\chronologyannotationclaim{composition.gospel-of-luke}{composition}{catholic-traditional-v1}{disputed}{before the end of the Roman imprisonment, when the Acts was finished}{Before Acts and the end of Paul’s Roman imprisonment}\\[.3em]%
}
\expandafter\def\csname triptychchronologyorientation@row-63f2e0b4f00a20c1\endcsname#1{%
Epistle: Rom. 13:11--14 & Composition: The Epistle to the Romans & \chronologyannotationreach{Rom.13.11}{Rom}{true}\chronologyannotationreach{Rom.13.12}{Rom}{true}\chronologyannotationreach{Rom.13.13}{Rom}{true}\chronologyannotationreach{Rom.13.14}{Rom}{true}\chronologyannotationclaim{composition.epistle-to-the-romans}{composition}{catholic-traditional-v1}{disputed}{(Romans), 57}{Proposals span A.D.~57--58}\chronologyannotationreach{Rom.13.11}{Rom}{true}\chronologyannotationreach{Rom.13.12}{Rom}{true}\chronologyannotationreach{Rom.13.13}{Rom}{true}\chronologyannotationreach{Rom.13.14}{Rom}{true}\chronologyannotationclaim{composition.epistle-to-the-romans}{composition}{catholic-traditional-v1}{disputed}{the Epistle to the Romans (58)}{}\\[.3em]%
}

% Generated. Do not edit.
% Complete orientation-v3 table; wording, dates and order belong to the corpus.
\begin{chronologysharedtable}
\chronologyorientationrow{row-ffde50926b618362}{}
\chronologyorientationrow{row-1884ebeacdf710a8}{}
\chronologyorientationrow{row-cdc5f95d42da0398}{}
\midrule
\chronologyorientationrow{row-e043e4e3e3d8610c}{}
\chronologyorientationrow{row-59e074ad4fb4bf82}{}
\chronologyorientationrow{row-8f3ee9c054245e05}{}
\midrule
\chronologyorientationrow{row-7d7d97f8358968c7}{}
\chronologyorientationrow{row-f5e722ff79a9e57e}{}
\chronologyorientationrow{row-6c4c7f6537f3262a}{}
\midrule
\chronologyorientationrow{row-09b31c90e5cc8331}{}
\chronologyorientationrow{row-cf755db1bdb2f73a}{}
\midrule
\chronologyorientationrow{row-eb87f9559d582d9f}{}
\chronologyorientationrow{row-16c5e1487f2aee6d}{}
\midrule
\chronologyorientationrow{row-03a893478d4fede6}{}
\chronologyorientationrow{row-ec61c0aa1cf6d38e}{}
\midrule
\chronologyorientationrow{row-63f2e0b4f00a20c1}{}
\end{chronologysharedtable}
{\footnotesize Psalm numbering: Vg. = Vulgate; Heb. = Hebrew.\par}


\subsection{Evidence for each date}
\label{sec:date-location-evidence}
{\small For each passage, as the Scriptural Date and Location table names it: what is dated, how the corpus weighs each answer, the source's own words, and the sources cited. A subject the corpus does not date is reported as such and given no date.\par}
{\small
\par\smallskip\noindent\textit{Introit: Ps. 24:1--3 (Vg.; Heb. 25), Ps. 24:4 (Vg.; Heb. 25)}\par\nobreak
\begin{itemize}[leftmargin=1.2em,itemsep=1pt,topsep=1pt,parsep=0pt]
\item \textbf{Traditional attribution} (disputed). David's reign in the usual chronology: B.C.~1055--1015 (Corbett also reports writers dating the reign 30--50 years later) (``reigned from 1055 to 1015 B.C.''; The Catholic Encyclopedia, Volume 4 (New York, 1908), article 'David, King,' printed p. 642, opening paragraph). Calendar orientation: c.~1055--1015 B.C. (usual chronology; Corbett also reports dates 30--50 years later) (The Catholic Encyclopedia, Volume 4 (New York, 1908), article 'David, King,' printed p. 642, opening paragraph). Place: Hebron, then Jerusalem (The Holy Bible, Douay-Rheims (Challoner revision), Project Gutenberg eBook 1581: 2 Kings 5:5).
\item \textbf{Historical setting}. No date: the chronology corpus has not yet dated this (research pending).
\item \textbf{Composition} (preferred). Psalms: Before c.~165 B.C. (``before the Maccabean period, around 165 B.C.; no individual psalm can be dated securely''; Catholic critical chronology; New American Bible, Revised Edition --- USCCB web state (2026), Psalms introduction).
\end{itemize}
\par\smallskip\noindent\textit{Offertory: Ps. 24:1--3 (Vg.; Heb. 25)}\par\nobreak
\begin{itemize}[leftmargin=1.2em,itemsep=1pt,topsep=1pt,parsep=0pt]
\item \textbf{Traditional attribution}. As under Introit: Ps. 24:1--3 (Vg.; Heb. 25), Ps. 24:4 (Vg.; Heb. 25).
\item \textbf{Historical setting}. As under Introit: Ps. 24:1--3 (Vg.; Heb. 25), Ps. 24:4 (Vg.; Heb. 25).
\item \textbf{Composition}. As under Introit: Ps. 24:1--3 (Vg.; Heb. 25), Ps. 24:4 (Vg.; Heb. 25).
\end{itemize}
\par\smallskip\noindent\textit{Gradual: Ps. 24:3 (Vg.; Heb. 25), Ps. 24:4 (Vg.; Heb. 25)}\par\nobreak
\begin{itemize}[leftmargin=1.2em,itemsep=1pt,topsep=1pt,parsep=0pt]
\item \textbf{Traditional attribution}. As under Introit: Ps. 24:1--3 (Vg.; Heb. 25), Ps. 24:4 (Vg.; Heb. 25).
\item \textbf{Historical setting}. As under Introit: Ps. 24:1--3 (Vg.; Heb. 25), Ps. 24:4 (Vg.; Heb. 25).
\item \textbf{Composition}. As under Introit: Ps. 24:1--3 (Vg.; Heb. 25), Ps. 24:4 (Vg.; Heb. 25).
\end{itemize}
\par\smallskip\noindent\textit{Alleluia: Ps. 84:8 (Vg.; Heb. 85)}\par\nobreak
\begin{itemize}[leftmargin=1.2em,itemsep=1pt,topsep=1pt,parsep=0pt]
\item \textbf{Historical setting}. As under Introit: Ps. 24:1--3 (Vg.; Heb. 25), Ps. 24:4 (Vg.; Heb. 25).
\item \textbf{Composition} (preferred). The Psalms of the Sons of Korah: Between the days of Isaias and the return from exile (The Catholic Encyclopedia, Volume 12 (1911)). Calendar orientation: c.~740--536 B.C. (Korahite collection estimate) (The Catholic Encyclopedia, Volume 12 (1911); The Catholic Encyclopedia, Volume 8 (New York, 1910); The Catholic Encyclopedia, Volume 3 (New York, 1908)).
\end{itemize}
\par\smallskip\noindent\textit{Communion: Ps. 84:13 (Vg.; Heb. 85)}\par\nobreak
\begin{itemize}[leftmargin=1.2em,itemsep=1pt,topsep=1pt,parsep=0pt]
\item \textbf{Historical setting}. As under Introit: Ps. 24:1--3 (Vg.; Heb. 25), Ps. 24:4 (Vg.; Heb. 25).
\item \textbf{Composition}. As under Alleluia: Ps. 84:8 (Vg.; Heb. 85).
\end{itemize}
\par\smallskip\noindent\textit{Gospel: Luke 21:25--33}\par\nobreak
\begin{itemize}[leftmargin=1.2em,itemsep=1pt,topsep=1pt,parsep=0pt]
\item \textbf{Narrated event} (preferred). Jesus foretells Jerusalem's fall and His coming: c.~A.D.~29 (derived; Maas's chronology) (The Holy Bible, Douay-Rheims (Challoner revision), Project Gutenberg eBook 1581: Luke 21:5, 21:7, 21:36, 21:37, 21:38; The Catholic Encyclopedia, Volume 8 (New York, 1910), article 'Chronology of the Life of Jesus Christ,' Ninth journey, heading and Holy Week paragraph: the close of the Tuesday; The Catholic Encyclopedia, Volume 8 (New York, 1910), A. J. Maas, 'Jesus Christ,' 'The Year of the Death of Christ,' printed p. 378, left column). Place: Jerusalem, by the Temple in Luke's narrative; Mount Olivet in Matthew and Mark (The Holy Bible, Douay-Rheims (Challoner revision), Project Gutenberg eBook 1581: Luke 21:5, 21:37, Matt. 24:3, Mark 13:3).
\item \textbf{Composition} (disputed). The Gospel of St Luke: c.~A.D.~70 (``About the year 70''; The Catholic Encyclopedia, Volume 14 (New York, 1912)); Before Acts and the end of Paul’s Roman imprisonment (``before the end of the Roman imprisonment, when the Acts was finished''; De auctore, de tempore compositionis et de historica veritate Evangeliorum secundum Marcum et secundum Lucam (1912), responsa-i-ix).
\end{itemize}
\par\smallskip\noindent\textit{Epistle: Rom. 13:11--14}\par\nobreak
\begin{itemize}[leftmargin=1.2em,itemsep=1pt,topsep=1pt,parsep=0pt]
\item \textbf{Composition} (disputed). The Epistle to the Romans: A.D.~57 (``(Romans), 57''; The Catholic Encyclopedia, Volume 11 (New York, 1911)); A.D.~58 (``the Epistle to the Romans (58)''; The Catholic Encyclopedia, Volume 5 (New York, 1909)).
\end{itemize}
\par}
